\chapter{拉普拉斯变换}

常规傅里叶变换要求函数绝对可积，对常数、三角函数、指数增长函数都没有普通意义；而工程中往往只关心 \(t\ge0\) 的一段过程，\(t<0\) 的情形不必知道。把"截断到 \(t\ge0\)"与"再乘一个指数衰减因子"两件事合起来，就得到拉普拉斯（Laplace）变换，简称拉氏变换。它是解线性常微分方程（组）与线性电路暂态过程最顺手的工具。

\section{拉普拉斯变换的定义与存在定理}

\subsection{定义}

设 \(f(t)\) 是定义在 \(t\ge0\) 上（含原点某邻域）的函数，\(s=\beta+i\omega\) 为复参数，定义

\keyformula{F(s)=\mathscr{L}[f(t)]=\int_0^{+\infty}f(t)e^{-st}\,dt}

称 \(F(s)\) 为 \(f(t)\) 的\textbf{像函数}，\(f(t)\) 为 \(F(s)\) 的\textbf{原函数}，记 \(f(t)=\mathscr{L}^{-1}[F(s)]\)。

\textbf{与傅里叶变换的关系}：在收敛域内取 \(s=\beta+i\omega\)，则
\[
F(s)=\mathscr{L}[f(t)]=\mathscr{F}\left[f(t)u(t)e^{-\beta t}\right](\omega).
\]
即"拉氏变换 = 先截断、再乘衰减因子、再作傅里叶变换"。指数因子 \(e^{-\beta t}\) 的作用正是把增长的函数"压"成绝对可积，这就是拉氏变换能处理 \(t^n\)、\(e^{at}\) 的原因。

由定义直接可得几个最基本的像函数：
\[
\mathscr{L}[1]=\mathscr{L}[u(t)]=\frac{1}{s}\ (\operatorname{Re}s>0),\qquad
\mathscr{L}\left[e^{at}\right]=\frac{1}{s-a}\ (\operatorname{Re}s>\operatorname{Re}a),
\]
\[
\mathscr{L}[\cos at]=\frac{s}{s^2+a^2},\qquad
\mathscr{L}[\sin at]=\frac{a}{s^2+a^2}\qquad(\operatorname{Re}s>0).
\]

\subsection{存在定理}

\textbf{定理}：设 \(f(t)\) 在 \(t\ge0\) 上连续（或分段连续），若存在正实数 \(c\) 与 \(M\)，使
\[
|f(t)|\le Me^{ct},\qquad t\ge0,
\]
则 \(F(s)=\mathscr{L}[f(t)]\) 在半平面 \(\operatorname{Re}s>c\) 上收敛（且在该半平面内解析）。

证明：由积分估值不等式
\[
\int_0^{+\infty}\left|f(t)e^{-st}\right|dt
\le M\int_0^{+\infty}e^{-(\operatorname{Re}s-c)t}dt
=\frac{M}{\operatorname{Re}s-c}<\infty\quad(\operatorname{Re}s>c).
\]

满足上述条件的函数称为\textbf{指数增长级函数}。\(t^n\)、\(e^{at}\)、\(\sin at\)、\(\cos at\) 及其乘积都是指数增长级，故都可作拉氏变换；而 \(e^{t^2}\)、\(\dfrac{1}{t}\)（在 \(t=0\) 附近无界过快）不满足条件，不能直接作拉氏变换。

\subsection{收敛域}

\textbf{定理}：设 \(f(t)\) 连续，\(F(s)=\mathscr{L}[f(t)]\)，则
\begin{enumerate}[label=(\arabic*)]
    \item 若积分在 \(s_0\) 处收敛，则它在半平面 \(\operatorname{Re}s>\operatorname{Re}s_0\) 上处处收敛；
    \item 若积分在 \(s_0\) 处发散，则它在半平面 \(\operatorname{Re}s<\operatorname{Re}s_0\) 上处处发散。
\end{enumerate}

证明思路（1）：构造 \(\Phi(t)=\displaystyle\int_0^t f(\tau)e^{-s_0\tau}d\tau\)，由假设 \(\lim_{t\to+\infty}\Phi(t)=F(s_0)\)，故 \(\Phi\) 有界；再由分部积分把 \(e^{-(s-s_0)t}\) 提到积分外，即得 \(\operatorname{Re}s>\operatorname{Re}s_0\) 时收敛。第 (2) 条是 (1) 的逆否命题。

由此，拉氏变换的收敛域一定是形如 \(\operatorname{Re}s>\alpha\) 的\textbf{半平面}（或全平面、或空集），边界 \(\operatorname{Re}s=\alpha\) 上的行为需单独讨论。收敛域与像函数一起才唯一确定原函数（因为不同的 \(f\) 在 \(t<0\) 的取值不影响像函数）。

\corepoint{三个固定认识：(1) 拉氏变换是\textbf{单边}变换，积分从 \(0\) 起，\(t<0\) 的信息被丢掉，所以逆变换得到的原函数默认"\(t<0\) 时取 0"；(2) \(s\) 是复参数，写 \(\operatorname{Re}s>c\) 而不是 \(s>c\)；(3) 本节以后一律默认 \(s\) 取在收敛域内，否则积分不收敛、像函数无意义。}

\section{拉普拉斯变换的性质}

以下设 \(\mathscr{L}[f(t)]=F(s)\)，\(\mathscr{L}[g(t)]=G(s)\)，且 \(s\) 都在相应收敛域内。

\subsection{线性性质}

对任意复数 \(\alpha,\beta\)，
\[
\mathscr{L}[\alpha f(t)+\beta g(t)]=\alpha F(s)+\beta G(s),\qquad
\mathscr{L}^{-1}[\alpha F(s)+\beta G(s)]=\alpha f(t)+\beta g(t).
\]

\subsection{相似、延迟与位移}

\textbf{相似性质}（\(a>0\)）：
\[
\mathscr{L}[f(at)]=\frac{1}{a}F\!\left(\frac{s}{a}\right).
\]

\textbf{延迟性质}（\(\tau>0\)）：
\[
\mathscr{L}\left[f(t-\tau)u(t-\tau)\right]=e^{-s\tau}F(s).
\]

\textbf{位移性质}（\(a\in\mathbb{C}\)）：
\[
\mathscr{L}\left[e^{at}f(t)\right]=F(s-a).
\]

\corepoint{延迟性质里被延迟的部分必须写成 \(f(t-\tau)u(t-\tau)\)：若只写 \(f(t-\tau)\)，则 \(t<\tau\) 时它并不为零，与"函数在 \(t<0\) 处为 0"的单边约定冲突，结果会多出错误项。另外注意"延迟"改的是像函数乘 \(e^{-s\tau}\)，"位移"改的是像函数的自变量 \(s-a\)，两者不可混。}

\subsection{微分性质}

\textbf{原函数求导}：
\[
\mathscr{L}\left[f'(t)\right]=sF(s)-f(0),
\]
\[
\mathscr{L}\left[f^{(n)}(t)\right]=s^nF(s)-s^{n-1}f(0)-s^{n-2}f'(0)-\cdots-f^{(n-1)}(0).
\]

\keyformula{\mathscr{L}\left[f^{(n)}(t)\right]=s^nF(s)-\sum_{k=1}^{n}s^{n-k}f^{(k-1)}(0)}

要求 \(f^{(k)}(t)\ (k\le n)\) 都是指数增长级。这组公式是"用拉氏变换解微分方程"的支点：求导变成乘 \(s\)，初始条件自动进入方程。

\textbf{像函数求导}：
\[
F'(s)=-\mathscr{L}[tf(t)],\qquad
F^{(n)}(s)=(-1)^n\mathscr{L}[t^nf(t)].
\]

\subsection{积分性质}

\textbf{原函数积分}：
\[
\mathscr{L}\left[\int_0^tf(\tau)\,d\tau\right]=\frac{F(s)}{s},
\qquad
\mathscr{L}\left[\underbrace{\int_0^t\!dt\int_0^t\!dt\cdots\int_0^t}_{n\ \text{次}}f(t)dt\right]=\frac{F(s)}{s^n}.
\]

\textbf{像函数积分}：
\[
\int_s^{+\infty}F(\sigma)\,d\sigma=\mathscr{L}\left[\frac{f(t)}{t}\right],
\]
\[
\underbrace{\int_s^\infty\!ds\int_s^\infty\!ds\cdots\int_s^\infty}_{n\ \text{次}}F(s)\,ds
=\mathscr{L}\left[\frac{f(t)}{t^n}\right].
\]

\subsection{卷积与卷积定理}

\textbf{拉氏卷积}（与傅里叶卷积不同：积分上限是 \(t\)，且两个因子都只取 \(t\ge0\) 的部分）：
\[
(f_1*f_2)(t)=\int_0^tf_1(\tau)f_2(t-\tau)\,d\tau
=\int_0^tf_1(t-\tau)f_2(\tau)\,d\tau
\qquad(t\ge0).
\]

\textbf{卷积定理}：
\[
\mathscr{L}[f_1*f_2](s)=F_1(s)F_2(s),
\qquad
\mathscr{L}^{-1}[F_1(s)F_2(s)]=f_1(t)*f_2(t).
\]
即像函数的乘积对应原函数的卷积。这与傅里叶变换的卷积定理完全平行，但拉氏卷积不必担心绝对可积性，工程上更常用。

\subsection{初值定理与终值定理}

\textbf{初值定理}：若 \(\lim_{s\to\infty}sF(s)\) 存在有限，则
\[
f(0^+)=\lim_{s\to+\infty}sF(s).
\]

\textbf{终值定理}：若 \(sF(s)\) 的奇点都位于左半平面（\(\operatorname{Re}s<0\)），且 \(\lim_{s\to0}sF(s)\) 存在有限，则
\[
f(+\infty)=\lim_{s\to0}sF(s).
\]

两个定理的价值在于：不必求逆变换，就能由像函数直接读出原函数的初值与终值（稳定性判断常用）。

\subsection{周期函数的像函数}

设 \(f(t)\) 是以 \(T\) 为周期的函数（\(t\ge0\)），则
\[
\mathscr{L}[f(t)]=\frac{1}{1-e^{-sT}}\int_0^Tf(t)e^{-st}\,dt .
\]
若 \(f(t)\) 是\textbf{交错周期函数}（\(f(t+T)=-f(t)\)），则
\[
\mathscr{L}[f(t)]=\frac{1}{1+e^{-sT}}\int_0^Tf(t)e^{-st}\,dt .
\]

\subsection{性质一览}

\begin{center}
\begin{tabular}{@{}lll@{}}
\toprule
性质 & 时域 & 像函数 \\
\midrule
线性 & \(\alpha f+\beta g\) & \(\alpha F+\beta G\) \\
相似 & \(f(at)\ (a>0)\) & \(\dfrac{1}{a}F\!\left(\dfrac{s}{a}\right)\) \\
延迟 & \(f(t-\tau)u(t-\tau)\) & \(e^{-s\tau}F(s)\) \\
位移 & \(e^{at}f(t)\) & \(F(s-a)\) \\
微分 & \(f^{(n)}(t)\) & \(s^nF(s)-\sum\limits_{k=1}^{n}s^{n-k}f^{(k-1)}(0)\) \\
像函数微分 & \(t^nf(t)\) & \((-1)^nF^{(n)}(s)\) \\
积分 & \(\displaystyle\int_0^tf(\tau)d\tau\) & \(\dfrac{F(s)}{s}\) \\
像函数积分 & \(\dfrac{f(t)}{t}\) & \(\displaystyle\int_s^{+\infty}F(\sigma)d\sigma\) \\
卷积 & \((f_1*f_2)(t)\) & \(F_1(s)F_2(s)\) \\
初值 & \(f(0^+)\) & \(\lim\limits_{s\to\infty}sF(s)\) \\
终值 & \(f(+\infty)\) & \(\lim\limits_{s\to0}sF(s)\) \\
\bottomrule
\end{tabular}
\end{center}

\section{拉普拉斯逆变换}

\subsection{反演积分与留数法}

由复反演公式，\(F(s)\) 的原函数可写成沿竖直直线的复积分
\[
f(t)=\frac{1}{2\pi i}\int_{\beta-i\infty}^{\beta+i\infty}F(s)e^{st}\,ds
\qquad(t>0),
\]
其中 \(\beta\) 取在收敛域内（即 \(F(s)\) 的所有奇点都在直线左侧）。

\textbf{定理}（留数法）：设 \(F(s)\) 在复平面上除有限个孤立奇点 \(s_1,s_2,\dots,s_n\) 外解析，且
\[
\lim_{s\to\infty}F(s)=0,
\]
则
\begin{equation}
f(t)=\mathscr{L}^{-1}[F(s)]=\sum_{k=1}^{n}\operatorname{Res}\left[F(s)e^{st},s_k\right]
\qquad(t>0),
\label{eq:invres}
\end{equation}
而 \(t<0\) 时 \(f(t)=0\)。

\textbf{计算步骤}：
\begin{enumerate}
    \item 验证 \(\lim_{s\to\infty}F(s)=0\)（分母次数高于分子次数即可），并找出 \(F\) 的全部奇点；
    \item 对 \(F(s)e^{st}\) 求各奇点的留数（一阶极点用 \(\lim(s-s_k)F(s)e^{st}\)；\(m\) 阶极点用求导公式）；
    \item 相加即得 \(t>0\) 时的 \(f(t)\)。
\end{enumerate}

\subsection{部分分式法}

把 \(F(s)\) 分解成若干"查表即得"的简单分式之和，是最常用的手段。设 \(F(s)=\dfrac{P(s)}{Q(s)}\)。

\begin{enumerate}
    \item \textbf{先化真分式}：若 \(\deg P\ge\deg Q\)，先作多项式除法，把多项式部分（对应 \(\delta\) 及其导数）与真分式分开；
    \item \textbf{分解分母}：\(Q(s)\) 分解为一次因式与不可约二次因式之积；
    \item \textbf{单极点}：\(\dfrac{A_k}{s-s_k}\)，其中
    \[
    A_k=\lim_{s\to s_k}(s-s_k)F(s)=\frac{P(s_k)}{Q'(s_k)}\quad(\text{若}\ s_k\ \text{是一阶零点});
    \]
    \item \textbf{\(m\) 重极点}：写成 \(\displaystyle\sum_{j=1}^{m}\frac{B_j}{(s-s_0)^j}\)，系数用
    \[
    B_{m-j+1}=\frac{1}{(j-1)!}\lim_{s\to s_0}\frac{d^{\,j-1}}{ds^{\,j-1}}\left[(s-s_0)^mF(s)\right]
    \]
    或直接用待定系数法；
    \item \textbf{共轭复极点}：把 \(\dfrac{Bs+C}{(s+\alpha)^2+\beta^2}\) 配成
    \[
    B\cdot\frac{s+\alpha}{(s+\alpha)^2+\beta^2}+\frac{C-B\alpha}{\beta}\cdot\frac{\beta}{(s+\alpha)^2+\beta^2},
    \]
    逆变换分别是 \(Be^{-\alpha t}\cos\beta t\) 与 \(\dfrac{C-B\alpha}{\beta}e^{-\alpha t}\sin\beta t\)。
\end{enumerate}

\subsection{卷积法}

当 \(F(s)=F_1(s)F_2(s)\) 且两个因子各自的逆变换容易求出时，用卷积定理
\[
f(t)=\mathscr{L}^{-1}[F_1F_2]=f_1(t)*f_2(t)=\int_0^tf_1(\tau)f_2(t-\tau)\,d\tau
\]
求原函数。它把"分解有理分式"的代数工作换成了"算一个卷积积分"，对含 \(e^{-s\tau}\) 的延迟型像函数尤其方便。

\corepoint{逆变换的三个易错点：(1) 留数法求的是 \(F(s)e^{st}\) 的留数，漏掉 \(e^{st}\) 就退回成求 \(F\) 的留数，结果必错；(2) 留数法要求 \(\lim_{s\to\infty}F(s)=0\)，对 \(F(s)=1\)（对应 \(\delta(t)\)）或多项式像函数不适用，必须另作处理；(3) 部分分式法必须先化成真分式，假分式的多项式部分对应的是 \(\delta(t)\) 及其导数，不能当作普通分式查表。}

\section{常用拉普拉斯变换对}

下表按教材「附录 II 拉氏变换简表」整理，约定 \(t\ge0\)（\(t<0\) 时各原函数取 0），\(m>-1\)，\(\Gamma\) 为伽马函数。

\begin{center}
\begin{tabular}{@{}ll@{}}
\toprule
\(f(t)\) & \(F(s)=\mathscr{L}[f(t)]\) \\
\midrule
\(1\ (u(t))\) & \(\dfrac{1}{s}\) \\
\(t\) & \(\dfrac{1}{s^2}\) \\
\(t^n\ (n=1,2,\dots)\) & \(\dfrac{n!}{s^{n+1}}\) \\
\(t^m\ (m>-1)\) & \(\dfrac{\Gamma(m+1)}{s^{m+1}}\) \\
\(\sqrt{t}\) & \(\dfrac{\sqrt{\pi}}{2s^{3/2}}\) \\
\(\dfrac{1}{\sqrt{t}}\) & \(\sqrt{\dfrac{\pi}{s}}\) \\
\(e^{at}\) & \(\dfrac{1}{s-a}\) \\
\(t^me^{at}\) & \(\dfrac{\Gamma(m+1)}{(s-a)^{m+1}}\) \\
\(\sin at\) & \(\dfrac{a}{s^2+a^2}\) \\
\(\cos at\) & \(\dfrac{s}{s^2+a^2}\) \\
\(\sinh at\) & \(\dfrac{a}{s^2-a^2}\) \\
\(\cosh at\) & \(\dfrac{s}{s^2-a^2}\) \\
\(e^{at}\sin bt\) & \(\dfrac{b}{(s-a)^2+b^2}\) \\
\(e^{at}\cos bt\) & \(\dfrac{s-a}{(s-a)^2+b^2}\) \\
\(t\sin at\) & \(\dfrac{2as}{(s^2+a^2)^2}\) \\
\(t\cos at\) & \(\dfrac{s^2-a^2}{(s^2+a^2)^2}\) \\
\(\dfrac{1}{a^2}\left(1-\cos at\right)\) & \(\dfrac{1}{s(s^2+a^2)}\) \\
\(\dfrac{1}{a^3}\left(at-\sin at\right)\) & \(\dfrac{1}{s^2(s^2+a^2)}\) \\
\(\dfrac{1}{a}\sin at-\dfrac{1}{b}\sin bt\) & \(\dfrac{b^2-a^2}{(s^2+a^2)(s^2+b^2)}\) \\
\(\cos at-\cos bt\) & \(\dfrac{(b^2-a^2)s}{(s^2+a^2)(s^2+b^2)}\) \\
\(\delta(t)\) & \(1\) \\
\(\delta^{(n)}(t)\) & \(s^n\) \\
\(f(t-\tau)u(t-\tau)\) & \(e^{-s\tau}F(s)\) \\
\(f(t)\ (T\ \text{为周期})\) & \(\dfrac{1}{1-e^{-sT}}\displaystyle\int_0^Tf(t)e^{-st}dt\) \\
\bottomrule
\end{tabular}
\end{center}

\section{拉普拉斯变换在解方程与电路中的应用}

\subsection{解线性常微分方程（组）}

利用微分性质，拉氏变换把"带初始条件的线性常微分方程"化成"像函数的代数方程"，解出像函数再逆变换即可。整个流程不需要先求通解、再代初始条件。

\textbf{步骤}：
\begin{enumerate}
    \item 设 \(\mathscr{L}[y(t)]=Y(s)\)，对方程两边作拉氏变换；
    \item 用微分性质把 \(y^{(k)}\) 换成 \(s^kY(s)-s^{k-1}y(0)-\cdots-y^{(k-1)}(0)\)，初始条件在这一步自动代入；
    \item 解出 \(Y(s)\)（有理分式）；
    \item 部分分式分解（或留数法）作逆变换，得 \(y(t)\)；
    \item 回代检验初值与终值（用初值、终值定理快速自查）。
\end{enumerate}

对\textbf{方程组}：设 \(\mathscr{L}[x]=X(s)\)、\(\mathscr{L}[y]=Y(s)\)，同样把每个方程变成 \(X,Y\) 的线性代数方程，解线性方程组后再逐个逆变换。这一步在电路（网孔法、节点法）里就是"把微分方程组变成复数代数方程组"。

\subsection{电路的 \(s\) 域模型}

对电路元件的时域关系作拉氏变换，得到\textbf{运算电路}（\(s\) 域模型）：

\begin{center}
\begin{tabular}{@{}lll@{}}
\toprule
元件 & 时域关系 & \(s\) 域关系 \\
\midrule
电阻 \(R\) & \(u=Ri\) & \(U(s)=RI(s)\) \\
电感 \(L\) & \(u=L\dfrac{di}{dt}\) & \(U(s)=sLI(s)-Li(0)\) \\
电容 \(C\) & \(i=C\dfrac{du}{dt}\) & \(U(s)=\dfrac{1}{sC}I(s)+\dfrac{u(0)}{s}\) \\
\bottomrule
\end{tabular}
\end{center}

即：电感等效为阻抗 \(sL\) 串联一个电压源 \(-Li(0)\)（或并联电流源 \(\dfrac{i(0)}{s}\)）；电容等效为阻抗 \(\dfrac{1}{sC}\) 串联电压源 \(\dfrac{u(0)}{s}\)（或并联电流源 \(Cu(0)\)）。初始条件由此变成"外加电源"，之后完全按直流电路的方法（串并联、节点法、网孔法）列代数方程。

\textbf{传递函数与响应}：
\begin{itemize}
    \item 令所有初始值为零，定义\textbf{传递函数}（网络函数）
    \[
    H(s)=\frac{Y(s)}{X(s)},
    \]
    它只由电路结构与参数决定，与输入无关；
    \item \textbf{冲激响应} \(h(t)=\mathscr{L}^{-1}[H(s)]\)，则任意输入下的零状态响应为 \(y(t)=h(t)*x(t)\)，即 \(Y(s)=H(s)X(s)\)；
    \item 全响应可分解为\textbf{零输入响应}（初始条件引起，令 \(X=0\)）与\textbf{零状态响应}（输入引起，令初始值为 0）之和；
    \item \(H(s)\) 的极点位置决定暂态衰减快慢与稳定性：极点全在左半平面（\(\operatorname{Re}s<0\)）时系统稳定，这与终值定理的条件一致。
\end{itemize}

\corepoint{解方程与电路应用中的高频错误：(1) 作变换时漏掉初始条件项（\(-f(0)\)、\(-s f(0)-f'(0)\) 等），这是最常见的失分点；(2) 逆变换前忘记把 \(Y(s)\) 化成真分式；(3) 混淆"零输入"与"零状态"——前者只由初始储能决定，后者只由外加激励决定，叠加时初始条件不能重复计入；(4) 用终值定理判稳态值时，忘记先确认 \(sF(s)\) 的极点全在左半平面（含虚轴极点时会得出错误结论）。}
